\part{Riferimento}

\chapter{Preset Comfort e Visibilità}
\label{ch:preset-comfort}

Usa sempre il preset più adatto alle condizioni di luce attuali.
Se un preset sembra attraente ma rende più difficile vedere le informazioni urgenti,
passa a uno più chiaro.

\section{Guida Pratica}

\begin{itemize}
  \item Usa preset più luminosi e ad alto contrasto alla luce del giorno.
  \item Usa preset più scuri orientati alla notte quando preservare la visione notturna
        è importante.
  \item Mantieni allarmi e testo critico per la sicurezza facilmente leggibili
        su ogni preset.
  \item Il preset Notte imposta l'icona tinta in rosso per impostazione predefinita.
        Questo è intenzionale --- la luce rossa minimizza il tempo necessario per
        riadattarsi all'oscurità dopo aver guardato il display.
\end{itemize}

\chapter{Display, Tocco e Operatività in Condizioni Difficili}
\label{ch:condizioni-difficili}

\section{Leggibilità del Display}

FairWindSK deve rimanere leggibile alla luce del sole piena, al crepuscolo, di notte
e in condizioni difficili. Inizia sempre con:

\begin{enumerate}
  \item Luminosità dello schermo.
  \item Preset comfort (corretto per le condizioni di luce attuali).
  \item Angolo di visione e posizione di plancia.
  \item L'app o la pagina mostrata nell'Area Applicazioni.
\end{enumerate}

\section{Uso del Tocco in Condizioni Difficili}

\begin{itemize}
  \item Prima mettiti in sicurezza. Una mano all'imbarcazione, una allo schermo.
  \item Usa tocchi deliberati e singoli. La pressione esitante e i tocchi rapidi
        ripetuti causano azioni indesiderate.
  \item Non esplorare mai i controlli dell'autopilota toccando in modo speculativo ---
        ogni comando ha un effetto fisico immediato.
  \item Mantieni lo schermo pulito. La salsedine riduce significativamente la precisione
        del tocco.
  \item Pre-configura quello che ti serve mentre le condizioni sono calme.
  \item Preferisci i target grandi e stabili delle icone principali della Barra Inferiore.
\end{itemize}

\chapter{Procedura Operativa Consigliata}
\label{ch:procedura}

\section{Prima della Partenza}

\begin{enumerate}
  \item Avvia FairWindSK e conferma che lo stato della shell sia normale.
  \item Verifica che i dati chiave attesi siano live nella Barra Superiore.
  \item Apri l'app più rilevante per la fase di navigazione corrente.
  \item Conferma luminosità e preset comfort.
  \item Completa la lista di controllo pre-partenza (Cap.~2).
\end{enumerate}

\section{Durante la Traversata}

\begin{enumerate}
  \item Controlla i valori obsoleti o mancanti nella Barra Superiore.
  \item Verifica rotta, profondità e prua quando le condizioni cambiano.
  \item Usa il launcher per cambiare attività invece di forzare un'app bloccata.
  \item Tratta gli stati degradati come avvisi operativi.
  \item Quando la schermata corrente smette di sembrare chiara o affidabile, tocca
        \iconlg{applications}~\key{App} per tornare al launcher.
\end{enumerate}

\section{Dopo una Riconnessione o Riavvio del Server}

\begin{enumerate}
  \item Attendi brevemente il completamento del recupero automatico.
  \item Conferma che lo stato mostri \textit{Live}.
  \item Riapri l'app se l'applicazione in primo piano era degradata.
  \item Riconferma i valori di navigazione chiave prima di affidarti ad essi.
\end{enumerate}

\chapter{Signal K Server su Docker}
\label{ch:docker}

\section{Immagine Docker Ufficiale}

\begin{lstlisting}
cr.signalk.io/signalk/signalk-server:latest
\end{lstlisting}

Architetture supportate: \code{linux/amd64}, \code{linux/arm/v7}, \code{linux/arm64}
(tutte le configurazioni Raspberry~Pi~4 e~5).

\section{Avvio Rapido con Dati NMEA Demo}

\begin{lstlisting}
mkdir $HOME/.signalk
docker run --init -it --rm \
  --name signalk-server \
  --publish 3000:3000 \
  -v $HOME/.signalk:/home/node/.signalk \
  cr.signalk.io/signalk/signalk-server \
  --sample-nmea0183-data
\end{lstlisting}

\section{Deployment in Produzione}

\begin{lstlisting}
docker run -d --init \
  --name signalk-server \
  -p 3000:3000 \
  -v $HOME/.signalk:/home/node/.signalk \
  cr.signalk.io/signalk/signalk-server
\end{lstlisting}

\section{Aggiornamento di Signal K}

\begin{lstlisting}
docker pull cr.signalk.io/signalk/signalk-server:latest
docker stop signalk-server
docker rm signalk-server
# Riesegui il comando docker run di produzione sopra
\end{lstlisting}

\chapter{Risoluzione dei Problemi}
\label{ch:problemi}

\section{Valori Mancanti}

\begin{enumerate}
  \item Controlla lo stato della shell.
  \item Verifica se altri valori sono ancora live.
  \item Conferma che lo strumento sorgente funzioni.
  \item Controlla \ui{Impostazioni~> Percorsi Signal~K} --- è configurato il percorso corretto?
  \item Torna al launcher e riapri l'app interessata se necessario.
\end{enumerate}

\section{Valori Non Aggiornati}

\begin{enumerate}
  \item Tratta i valori con cautela.
  \item Attendi brevemente il recupero.
  \item Verifica con altri strumenti.
  \item Indaga se lo stato persiste oltre 60 secondi.
\end{enumerate}

\section{Un'App Web è Vuota o Bloccata}

\begin{enumerate}
  \item Riprova una volta se la shell lo propone.
  \item Torna al launcher.
  \item Controlla la salute della shell e lo stato della connessione Signal~K.
  \item Usa un'altra app se la shell è altrimenti integra.
  \item Un'app in errore non significa che FairWindSK abbia smesso di funzionare.
\end{enumerate}

\section{Shell Disconnessa}

\begin{enumerate}
  \item Signal~K è in esecuzione? Controlla i LED sul Raspberry~Pi; controlla lo stato
        del container Docker: \code{docker ps -a | grep signalk}.
  \item Il dispositivo è sulla rete Wi-Fi corretta?
  \item L'indirizzo del server in \ui{Impostazioni~> Connessione} è corretto?
  \item Un browser su un altro dispositivo raggiunge \code{http://indirizzo-server:3000}?
  \item Attendi 60 secondi per il tentativo di riconnessione automatica.
  \item Riconferma tutti i valori critici per la sicurezza dopo la riconnessione.
\end{enumerate}

\section{FairWindSK Lento su Raspberry Pi}

\begin{enumerate}
  \item Controlla la temperatura CPU: \code{vcgencmd measure\_temp}.
        Sopra 80°C indica throttling termico --- migliora il raffreddamento.
  \item Controlla \ui{Impostazioni~> Sistema~> Prestazioni} per l'utilizzo della memoria.
  \item Usa una scheda microSD Classe~A2 ad alta velocità.
        Una scheda lenta è la causa più comune di prestazioni scadenti su Raspberry~Pi.
  \item Imposta il Livello di Log su \textit{Nessun log} in
        \ui{Impostazioni~> Sistema~> Diagnostica}.
\end{enumerate}

\chapter{Lista di Controllo Operativa Giornaliera}
\label{ch:checklist}

\section*{Prima della Partenza}

\rowcolors{2}{LtGray}{white}
\begin{longtable}{C{0.8cm}L{12.9cm}}
\toprule
\rowcolor{Navy}
\color{white}\sffamily\bfseries $\square$ &
\color{white}\sffamily\bfseries Controllo \\
\midrule
$\square$ & FairWindSK avviato --- stato mostra Live \\
$\square$ & Posizione live e plausibile \\
$\square$ & SOG in aggiornamento --- fix GPS valido \\
$\square$ & Profondità in aggiornamento (se installata) \\
$\square$ & Rotta (HDG) visualizzata \\
$\square$ & Plotter aperto; imbarcazione nella posizione corretta \\
$\square$ & Rotta o waypoint attivo (se si procede verso una destinazione) \\
$\square$ & Target AIS visibili (se AIS installato) \\
$\square$ & Preset comfort impostato per le condizioni di luce attuali \\
$\square$ & Luminosità schermo leggibile dalla plancia \\
$\square$ & Posizione del pulsante POB confermata \\
$\square$ & Autopilota disponibile e funzionante (se richiesto) \\
$\square$ & Guardia ancora disponibile (se si prevede di ancorare) \\
\bottomrule
\end{longtable}

\section*{Dopo l'Arrivo}

\rowcolors{2}{LtGray}{white}
\begin{longtable}{C{0.8cm}L{12.9cm}}
\toprule
\rowcolor{Navy}
\color{white}\sffamily\bfseries $\square$ &
\color{white}\sffamily\bfseries Controllo \\
\midrule
$\square$ & Torna al launcher --- nessun allarme attivo \\
$\square$ & Guardia ancora impostata con raggio appropriato (se all'ancora) \\
$\square$ & Preset comfort cambiato in Crepuscolo o Notte \\
$\square$ & Traccia esportata per il giornale di bordo (se richiesto) \\
\bottomrule
\end{longtable}

\chapter{Scheda di Riferimento Rapido}
\label{ch:riferimento-rapido}

\section*{Guida alle Azioni Immediate}

\rowcolors{2}{LtGray}{white}
\begin{longtable}{L{3.6cm}L{4.4cm}L{5.7cm}}
\toprule
\rowcolor{Navy}
\color{white}\sffamily\bfseries Situazione &
\color{white}\sffamily\bfseries Azione immediata &
\color{white}\sffamily\bfseries Poi \\
\midrule
Persona in mare &
  Tocca \iconlg{alarm-pob}~\key{POB}. Una volta. &
  Chiama MAYDAY sul VHF Can.~16. Usa percorso SAR. \\
Allarme lampeggiante &
  Tocca \iconlg{alarm-general}~\key{Allarmi}. Leggi. &
  Agisci secondo le procedure di emergenza. \\
Schermo bloccato &
  Attendi 10~s. Tocca \iconlg{applications}~\key{App}. &
  Controlla lo stato. Prova a riaprire l'app. \\
Shell disconnessa &
  Attendi 60~s per la riconnessione automatica. &
  Se ancora disconnessa: controlla Signal~K e Wi-Fi. \\
Profondità obsoleta in acque basse &
  Verifica subito con uno strumento di backup. &
  Riduci la velocità. Usa scandaglio a mano o ecoscand. di backup. \\
Autopilota non risponde &
  Tocca \key{Standby}. Prendi il timone. &
  Indaga la causa prima di innestare di nuovo. \\
\bottomrule
\end{longtable}

\section*{Navigazione con Un Tocco}

\rowcolors{2}{LtGray}{white}
\begin{longtable}{L{5.8cm}L{7.9cm}}
\toprule
\rowcolor{Navy}
\color{white}\sffamily\bfseries Obiettivo &
\color{white}\sffamily\bfseries Azione \\
\midrule
Torna alla schermata principale  & Tocca \iconlg{applications}~\key{App} --- sempre disponibile \\
Passa a un'altra app aperta      & Tocca la sua icona nella striscia app aperte \\
Cambia preset comfort            & \ui{Impostazioni~> Comfort} \\
Controlla lo stato connessione   & Area di stato nella Barra Superiore (destra) \\
Segna posizione ancora           & Tocca \iconlg{anchor-iec}~\key{Ancora} \textrightarrow{} \iconlg{anchor-drop}~\key{Cala} \\
Annulla guardia ancora           & Tocca \iconlg{anchor-iec}~\key{Ancora} \textrightarrow{} \iconlg{anchor-raise}~\key{Salpa} \\
Annulla un evento POB            & Barra POB \textrightarrow{} seleziona incidente \textrightarrow{} \key{Annulla} \\
Esporta traccia traversata       & \ui{MyData~> Tracce} \textrightarrow{} \key{Esporta} \\
\bottomrule
\end{longtable}

\chapter{Note di Sicurezza}
\label{ch:sicurezza}

\begin{warningbox}
FairWindSK è un ausilio alla navigazione.
Non sostituisce la perizia nautica, la guardia visiva, le carte nautiche cartacee
o qualsiasi procedura di sicurezza certificata.
Il comandante e l'equipaggio dell'imbarcazione sono i soli responsabili della condotta
in sicurezza del viaggio.
\end{warningbox}

\begin{warningbox}
Verifica sempre le informazioni critiche per la sicurezza --- posizione, profondità, rotta,
rischio di collisione --- con strumenti indipendenti prima di agire su di esse.
Se FairWindSK contraddice ciò che vedi fuori dall'imbarcazione o su strumenti affidabili,
dai priorità all'osservazione diretta e ai riferimenti di navigazione indipendenti.
\end{warningbox}

\begin{warningbox}
Profondità mancante non è profondità zero.
Un valore mancante significa che la misura non è disponibile.
Trattala come completamente sconosciuta e naviga con la massima cautela.
\end{warningbox}

\begin{warningbox}
I comandi all'autopilota devono essere confermati rispetto al comportamento reale
dell'imbarcazione.
Non dare mai per scontato che un cambio di modalità sia avvenuto perché il pulsante
ha risposto.
Sii pronto a disinnestare e prendere il timone manuale in qualsiasi momento.
\end{warningbox}

\begin{warningbox}
L'attivazione di un allarme in FairWindSK genera una notifica Signal~K solo sulla rete
di bordo.
Non contatta la guardia costiera, l'MRCC o alcuna autorità esterna.
Effettua una chiamata di soccorso DSC e/o vocale sul Canale VHF~16 in qualsiasi
emergenza reale.
\end{warningbox}

\begin{warningbox}
In un'emergenza, dai priorità alla gestione dell'imbarcazione, alla sicurezza
dell'equipaggio e all'osservazione diretta rispetto a qualsiasi interazione
con un display elettronico.
Uno schermo è uno strumento; la tua perizia nautica è la tua sicurezza.
\end{warningbox}

\vfill
\begin{center}
  {\color{Teal}\rule{0.5\textwidth}{1pt}}\\[10pt]
  {\Large\sffamily\bfseries\color{Navy}Buona navigazione.}\\[8pt]
  {\normalsize\sffamily\color{gray}
    Il team OpenFairWind \textbullet{}
    Università degli Studi di Napoli 'Parthenope'}\\[4pt]
  {\small\sffamily\color{Teal}\url{https://github.com/OpenFairWind/FairWindSK}}
\end{center}
